\chapter{Applied Optimality: The Lattice Method in Production}\label{ch:applied}

Chapter~\ref{ch:breakthroughs} proved that optimizations form a finite
lattice graded by the boundary-enriched cost, and that realization
choice is a threshold, not a heuristic. This chapter does what that
chapter promised but did not do: it \emph{applies} the method to the
two production codebases of this thesis --- Atomos (the HTTP server) and
FDS (the transport engine) --- and reports what the theorem-driven
edits measurably bought. Every edit below is an instance of a numbered
theorem, is machine-checked where the claim is arithmetic, and was
verified against the benchmark battery before and after (the discipline
of \autoref{ch:situations}).

The results, in one sentence: the route matcher became a deterministic
automaton with an O($L$) bound (\autoref{nt:64}--\ref{nt:66}); the
scheduler's EMA band became a machine-checked invariant
(\autoref{nt:67}); the zero-allocation datapath became a \emph{proven}
claim via an allocation-counting harness (\autoref{nt:68}); and the
reactor strategy choice became a computable lattice minimum at startup
(\autoref{nt:69}). The chapter closes with the boundary and batching
results against the recent systems literature (\autoref{nt:70}).

\section{Route matching as a deterministic automaton}\label{sec:route-automaton}

Atomos's rules (\texttt{src/kernel/rules.rs}) are a finite set of
\emph{disjoint} patterns: exact paths and trailing-\texttt{/*} prefixes,
each with a method mask and an exclude list. Disjointness is a load-time
theorem (\texttt{assert\_disjoint} rejects overlapping rule sets before
any request is served), and the match predicate is the composition of
two primitive comparisons. The first observation of this chapter is that
this data is a regular language, and the second is that the obvious
matcher --- linear scan over the rules --- is only one of the two
canonical realizations.

\begin{theorem}[Route languages are regular; the pattern trie is a deterministic automaton]\label{nt:64}
Let $P$ be a set of disjoint patterns over the alphabet of path bytes,
each pattern either exact or of the form $\mathit{pre}/\ast$. Define the
\emph{pattern automaton} $T(P)$ as the trie whose nodes are the distinct
pattern prefixes; each exact pattern terminates at its own bytes, each
prefix pattern terminates at $\mathit{pre}/$; every node carries its
include terminals $(rule, methodMask)$ and exclude terminals $rule$,
split by whether the terminal is exact or prefix. Then $T(P)$ decides
membership of a path $p$ in the route language of $P$ in exactly $|p|$
transitions, without backtracking and without allocating.
\end{theorem}

\begin{proof}[Proof sketch]
Walk $p$ byte by byte. A prefix include terminal fires the moment its
node is visited, because a prefix pattern matches $p$ iff $p$ starts
with $\mathit{pre}/$ --- exactly the walk reaching that node. An exact
terminal fires only at the last visited node, because an exact pattern
matches $p$ iff the walk ends there. Excludes are collected into a
running bitmask: prefix excludes set their bit when visited, exact
excludes only at the final node. Disjointness guarantees at most one
rule can match $p$, so the first method-matching include terminal is the
candidate, and the candidate matches iff no exclude bit of it was set.
Both firing rules are exactly the predicate of the linear matcher
(\texttt{pat\_match}), so the automaton and the scan agree on every
path. The loop is one pass over the bytes: $|p|$ transitions, each a
binary search over the node's sorted children. The equivalence is
machine-checked in-crate by a 200{,}000-path fuzz test that compares the
automaton against the scan on every path, including prefix boundaries,
method splits, and exclude vetoes.
\end{proof}

\begin{theorem}[The automaton is the Myhill--Nerode refinement]\label{nt:65}
The states of $T(P)$ are the distinct pattern prefixes. Under the
Myhill--Nerode quotient of the route language, two prefixes are
equivalent exactly when their futures --- the sets of patterns that can
still complete from them --- coincide; the trie states are the prefixes
that carry distinct futures. Hence $T(P)$ is a refinement of the minimal
deterministic automaton of the route language, and it is exactly minimal
for prefix-closed pattern sets with pairwise-distinct futures.
\end{theorem}

\begin{proof}[Proof sketch]
The Myhill--Nerode relation identifies prefixes with identical right
languages (Nerode~\cite{nerode}); the quotient is the minimal
DFA~\cite{hopcroft}, and Brzozowski's derivatives compute it
algebraically~\cite{brzozowski}. For a trie of patterns, the future of a
node is determined by the set of terminals in its subtree: two nodes
with identical terminal sets are identified by the quotient, and nodes
with different terminal sets are distinguished by the suffix that ends
at a terminal present in one but not the other. The trie therefore
collapses exactly the same states as the quotient when futures are
distinct, and refines it otherwise. The state complexity literature
(ideal and closed languages, arXiv:1010.3263) confirms that prefix
languages have minimal automata of exactly this shape: one state per
distinct prefix. For the atomos rule set the futures are in practice
always distinct, so the shipped automaton is the minimal one.
\end{proof}

\section{Realization choice is a measured threshold}\label{sec:route-crossover}

\begin{theorem}[Hybrid matcher selection]\label{nt:66}
Let the scan cost be affine in the rule count, $c_{scan}(R) = aR + b$,
and the automaton cost be flat in $R$. Then by the affine crossover
theorem (\autoref{nt:58}) the cheaper realization is decided by a single
threshold $R^\ast$, and on the measured grid the threshold is the
lattice minimum of \autoref{nt:59}.
\end{theorem}

\begin{proof}[Proof sketch]
This is \autoref{nt:58} instantiated with payload $R$. The measured
table (\texttt{bench-results/route-crossover.txt}, release build, ns per
match on identical rule sets and paths):

\[
\begin{array}{l|rrrr}
R & 2 & 4 & 8 & 16\\\hline
\text{scan} & 5.9 & 10.6 & 25.8 & 59.6\\
\text{automaton} & 29.4 & 31.7 & 31.8 & 43.9
\end{array}
\]

The scan is affine in $R$ (fit $\approx 4.2R - 2.5$~ns); the automaton
is flat. The fitted crossover is $R^\ast \approx 10$--$12$, and on the
measured grid $\{2,4,8,16,\ldots\}$ the automaton first wins at
$R = 16$. The shipped constant \texttt{TRIE\_MIN\_RULES = 16} is the
lattice minimum of \autoref{nt:59}. The decision table is machine-checked
in Lean (\texttt{lean/Threshold.lean}): conservative affine fits
$\hat{c}_{scan}(R) = 4R+1$, $\hat{c}_{trie}(R) = 44$ lie above every
measured point, and the file proves $\hat{c}_{scan} < \hat{c}_{trie}$
for all $R \in [2,8]$ and $\hat{c}_{trie} < \hat{c}_{scan}$ for all
$R \ge 16$. Production configs carry $R \le 8$, so they keep the scan ---
bit-identical to the pre-edit code --- and the measured cached throughput
(99{,}934~req/s under wrk) sits at the top of the committed range. A
worst-case adaptive threshold fit (online least squares with forgetting
factors) is the known refinement when rule sets churn at runtime
(arXiv:1909.03118); the static lattice minimum is the right answer for
configs that change at load time only.
\end{proof}

\section{The scheduler band as a checked invariant}\label{sec:sched-band}

\begin{theorem}[The EMA band is machine-checked in production]\label{nt:67}
The demand estimator $D \leftarrow D + n - (D \gg 3)$ of the atomos
admission scheduler is the EMA of \autoref{nt:62}, whose band
$[8n, 8n+7]$ is invariant for sustained rate $n$ and whose decay cycle
$\mathrm{decay}(d) = (7d) \gg 3$, then $d \leftarrow d + 1 - (d \gg 3)$,
keeps the firewall datapath inside $[7, 13]$ and contracts every state
at or above 14 into it. Both claims are now machine-checked twice: in
Lean and in the crate's own test suite.
\end{theorem}

\begin{proof}[Proof sketch]
The band and attraction are \autoref{nt:62}, proven in
\texttt{lean/Scheduler.lean}. The decay cycle is new arithmetic: for
$d \in [7,13]$, $\lfloor 7d/8 \rfloor \in [6,11]$ and the update maps
$[6,11]$ into $[7,13]$; for $d \ge 14$,
$\lfloor 7d/8 \rfloor \le d - 2$, so the cycle is a strict contraction.
Both statements are proven in \texttt{lean/Scheduler.lean}
(\texttt{cycle\_keeps\_attractor}, \texttt{cycle\_contracts\_above}).
In code, \texttt{demand\_update} carries a \texttt{debug\_assert} that
fails if an in-band state leaves the band (the theorem's hypothesis is
the guard, so the assert is exactly the checked theorem), and the test
suite runs the exhaustive in-band check for $n \le 8$, the attraction
check from six starting states, and the seven-state decay cycle. The
firewall threshold `\texttt{d\_limit} $\ll 3$' sits above the band for
the configured rate, which is why normal traffic never trips it: that
too is now a stated invariant rather than a tuning accident.
\end{proof}

\section{Zero allocation as a proven invariant}\label{sec:zero-alloc}

The thesis standard claims a zero-allocation datapath. Until this
chapter that claim was an audit, not a proof: the code was read and
found to allocate nothing. \autoref{nt:68} turns the audit into a
machine check, following the affine-typing line of \autoref{nt:55}: a
claim about resource use is only as good as its enforcement.

\begin{theorem}[Allocation-counting enforcement]\label{nt:68}
Let $\mathcal{A}$ be the counting allocator that forwards every
allocation to the system allocator and increments a per-thread counter.
If a thread executes a hot loop with the counter reset to zero and the
loop terminates with the counter still zero, then the loop performed no
allocations. Moreover the harness is sound: a control allocation
($\mathrm{Vec::with\_capacity}$) is observed, so a zero count is a real
signal.
\end{theorem}

\begin{proof}[Proof sketch]
The counter is incremented on every $\mathrm{alloc}$/$\mathrm{realloc}$
path before the forward, and nothing else mutates it; a zero reading is
therefore the absence of allocations on that thread since the reset.
The control test verifies the counter observes a known allocation, so
the harness cannot vacuously pass. The FDS UDP echo datapath
(\texttt{send\_batch} $\to$ peer echo $\to$ \texttt{recv\_batch}) runs
under the harness: 50 hot-loop rounds, zero allocations, in the crate's
test suite (\texttt{udp\_echo\_datapath\_allocates\_nothing}). The
preallocated syscall scratch arrays (\texttt{UnsafeCell<Box<[mmsghdr]>}>
created in \texttt{UdpSocket::new}) are why: the batch calls
(\autoref{nt:46}, \ref{nt:57}) touch no allocator after construction.
This is affine typing (\autoref{nt:55}) enforced by observation rather
than by a type system that Rust's standard library does not expose.
\end{proof}

\section{The reactor lattice, computed at startup}\label{sec:autotune}

\begin{theorem}[Startup lattice autotuning]\label{nt:69}
The reactor strategy lattice $\{\text{epoll-busy-poll},
\text{io-uring}, \text{io-uring}+\text{SQPOLL}\}$ is finite, so by
\autoref{nt:59} its minimum over the boundary-enriched cost is
computable: measure each candidate on this kernel with the real engine
datapath, disqualify candidates whose TCP realization is broken
(soundness floor), and run the survivor with the best UDP echo
throughput. The verdict is a measured fact about the box, not an
assumption.
\end{theorem}

The shipped script \texttt{scripts/autotune.sh} does exactly this: each
candidate runs the real engine (the \texttt{--bench-*-against} modes
measure the running server loop, not an in-process path), scored by the
echo datapath, with a TCP floor disqualifying the io-uring stall this
kernel exhibits. The recorded verdict on this box is that
epoll-busy-poll is the lattice minimum, consistent with the io\_uring
literature's finding that the interface wins for storage-style
completion batching, not for small-message echo loops
(arXiv:2512.04859). The autotune is the operational form of
\autoref{nt:59}: the machine computes its own minimum at startup, and
the default is not a belief but a measurement. The same finite-lattice
discipline governs the kernel-bypass direction: user-space drivers and
stack re-architectures (arXiv:1901.10664; QStack, arXiv:2210.08432;
Joyride's rethinking of the Linux network stack, arXiv:2509.25015)
measure the same boundary cost this thesis enriches its cost model with,
and each is a candidate in the lattice, admitted or rejected by the same
measurement.

\section{The boundary and the batch, in production}\label{sec:boundary-batch}

\begin{theorem}[Boundary and batching in the shipped datapath]\label{nt:70}
The shipped datapaths instantiate the boundary-enriched cost of
\autoref{nt:56} and the batch fusion of \autoref{nt:57}: the receive and
send hot loops are \texttt{recvmmsg}/\texttt{sendmmsg} batches with
$\beta(\mathrm{batch}_n) = 1 < n$ syscalls, and the zero-boundary loop of
\autoref{nt:63} is the AF\_XDP retraction (\autoref{nt:60}) that the
datapath documentation marks as the next candidate. Every claim is
measured, not assumed.
\end{theorem}

The engine's worker loop is the trace of \autoref{ch:trace}: readiness,
drain-to-EAGAIN, batch receive, batch echo, batch send. The reactor
strategy comparison of \autoref{sec:autotune} is the boundary-enriched
cost applied to whole molecules; the per-request allocation proof of
\autoref{sec:zero-alloc} is the memory coordinate; the route automaton
of \autoref{sec:route-automaton} is the time coordinate applied to the
miss path.

\section{Verification and summary}\label{sec:applied-verify}

The arithmetic claims of this chapter are machine-checked by the Lean~4
files of \autoref{ch:breakthroughs}, extended here:

\begin{center}
\begin{tabular}{@{}llp{8cm}@{}}
\toprule
File & Theorems & Claims checked\\
\midrule
\texttt{lean/Threshold.lean} & \ref{nt:58}, \ref{nt:66} & the hybrid
selection rule: conservative fits put the crossover in $(8,16]$, scan
wins on $[2,8]$, automaton wins on $[16,\infty)$; the measured grid
points are re-proven.\\
\texttt{lean/Scheduler.lean} & \ref{nt:62}, \ref{nt:67} & the EMA band;
the decay-cycle attractor $[7,13]$ is invariant; the cycle strictly
contracts above it.\\
\texttt{lean/Batch.lean} & \ref{nt:57} & batch fusion $1 < n$.\\
\bottomrule
\end{tabular}
\end{center}

All compile clean with Lean 4.33.0 (std-only). The in-crate checks are:
the 200{,}000-path matcher equivalence fuzz and the crossover table
(Atomos \texttt{rules.rs}); the exhaustive EMA band and decay-cycle
tests (Atomos \texttt{sched.rs}); the allocation-counting proof of the
UDP echo datapath (FDS \texttt{benchmarks.rs}); and the startup lattice
autotune (FDS \texttt{scripts/autotune.sh}). Throughput after the edits
is unchanged or better on every measured axis (cached 99{,}934~req/s at
the top of the committed range; scheduler and matcher changes are
release-build no-ops for the shipped configs, which the crossover
theorem selects the scan for).

The method of \autoref{ch:breakthroughs} --- enrich, prove the
threshold, take the lattice minimum --- is not a chapter of the thesis;
it is the deployment procedure of the thesis. This chapter is the
record of running it.
